\needspace{8\baselineskip}
\subsection*{Inner products and their norms}

\begin{problem}\label{ex:06:examples}
Verify the inner-product axioms for (a) \(\R^m\) with \(\sum_i x_i y_i\);
(b) real \(m\times n\) matrices with \(\tr(A^\top B)\).
\end{problem}
\begin{solution}
In (a), symmetry and linearity follow term by term; \(\sum_i x_i^2\)
is nonnegative and vanishes precisely when every coordinate is zero.
In (b), \(\tr(A^\top B)=\sum_{i,j}a_{ij}b_{ij}\), so the same argument
applies to the entries.
\end{solution}

\begin{problem}\label{ex:06:schwarz}
In a real inner-product space prove \(|\ip{x}{y}|\leqslant\norm{x}\norm{y}\).
Determine exactly when equality holds, including zero vectors.
\end{problem}
\begin{solution}
If \(y=0\), both sides vanish. For \(y\ne0\), put
\(r=x-\ip{x}{y}y/\norm{y}^2\). Then
\[
0\leqslant\norm{r}^2=\norm{x}^2-\frac{\ip{x}{y}^2}{\norm{y}^2}.
\]
Multiplication by \(\norm{y}^2\) proves the inequality. Equality holds
exactly when \(r=0\), or \(x\) is a scalar multiple of \(y\).
Including the case \(y=0\), the condition is that the pair is linearly dependent.
\end{solution}

\begin{problem}\label{ex:06:parallelogram}
Prove \(\norm{x+y}^2+\norm{x-y}^2=2\norm{x}^2+2\norm{y}^2\)
and explain its name, the parallelogram identity.
\end{problem}
\begin{solution}
The two squares expand to \(\norm{x}^2\pm2\ip{x}{y}+\norm{y}^2\);
adding cancels the mixed terms. A parallelogram with adjacent side vectors
\(x,y\) has diagonals \(x+y,x-y\). Thus the sum of the squared diagonal
lengths equals the sum of the squared lengths of its four sides.
\end{solution}

\begin{problem}\label{ex:06:norm}
A norm satisfies \(\nu(ax)=|a|\nu(x)\), positivity with equality only at
zero, and \(\nu(x+y)\leqslant\nu(x)+\nu(y)\).
Prove that \(\nu(x)=\sqrt{\ip{x}{x}}\) is a norm. Include the Frobenius
norm as a special case.
\end{problem}
\begin{solution}
Bilinearity gives \(\ip{ax}{ax}=a^2\ip{x}{x}\), proving homogeneity.
Positive definiteness gives positivity and its equality condition.
Schwarz gives
\[
\norm{x+y}^2=\norm{x}^2+2\ip{x}{y}+\norm{y}^2
\leqslant(\norm{x}+\norm{y})^2.
\]
Taking nonnegative square roots proves the triangle inequality.
For matrices the product is \(\tr(A^\top B)\), so the norm is
\(\sqrt{\sum_{i,j}a_{ij}^2}\) and all three conclusions apply.
\end{solution}

\begin{problem}\label{ex:06:polarization}
Prove that a real inner product is recovered from its norm by
\[\ip{x}{y}=\tfrac14(\norm{x+y}^2-\norm{x-y}^2).\]
\end{problem}
\begin{solution}
Expanding the two squares and subtracting cancels the two pure-square
terms and leaves \(4\ip{x}{y}\). In particular two inner products
inducing the same norm must agree on every pair.
\end{solution}

\begin{problem}\label{ex:06:noninner-norm}
For \(m\geqslant2\), let \(\nu(x)=\sum_{i=1}^m|x_i|\) on \(\R^m\).
Prove that \(\nu\) is a norm but is not induced by an inner product.
\end{problem}
\begin{solution}
Absolute values give \(\nu(ax)=|a|\nu(x)\) and
\(\nu(x+y)\leqslant\sum_i(|x_i|+|y_i|)=\nu(x)+\nu(y)\).
Every summand is nonnegative; their sum is zero exactly when \(x=0\).
For distinct coordinate vectors \(e_1,e_2\), however,
\[
\nu(e_1+e_2)^2+\nu(e_1-e_2)^2=8,
\qquad 2\nu(e_1)^2+2\nu(e_2)^2=4.
\]
The parallelogram identity fails, so no inner product induces this norm.
\end{solution}

\begin{problem}\label{ex:06:distance}
For \(d(x,y)=\norm{x-y}\), prove (a) \(d(x,x)=0\);
(b) \(d(x,y)>0\) for \(x\ne y\); (c) symmetry;
(d) \(d(x,y)\leqslant d(x,z)+d(z,y)\). Interpret these in \(\R^2\).
\end{problem}
\begin{solution}
The first two assertions are norm positivity. Symmetry follows from
\(\norm{-(x-y)}=\norm{x-y}\). Finally
\(x-y=(x-z)+(z-y)\), so the norm triangle inequality proves (d).
In the plane, coincident points have distance zero and distinct points
have positive distance; reversing a route does not change its length;
and the direct segment is no longer than the two-segment route through \(z\).
\end{solution}

\begin{problem}\label{ex:06:continuity}
Suppose \(\norm{x_k-x}\to0\) and \(\norm{y_k-y}\to0\) in a real
inner-product space. Prove convergence of the norms and of
\(\ip{x_k}{y_k}\) to \(\ip{x}{y}\).
\end{problem}
\begin{solution}
The triangle inequality in both directions gives
\(\big|\norm{x_k}-\norm{x}\big|\leqslant\norm{x_k-x}\to0\),
and similarly for \(y_k\). In particular \(\norm{x_k}\) is bounded.
By bilinearity and Schwarz,
\[
|\ip{x_k}{y_k}-\ip{x}{y}|
\leqslant\norm{x_k}\norm{y_k-y}+\norm{x_k-x}\norm{y}\longrightarrow0.
\]
\end{solution}

\needspace{8\baselineskip}
\subsection*{Orthogonality and construction}

\begin{problem}\label{ex:06:orthogonal-vectors}
Show that zero is orthogonal to every vector. Find every vector in
\(\R^3\) orthogonal to \((2,2,-1)^\top\).
\end{problem}
\begin{solution}
Linearity gives \(\ip{0}{v}=\ip{0+0}{v}=2\ip{0}{v}\), hence zero.
The coordinate condition is \(2x+2y-z=0\); its complete solution is
\((s,t,2s+2t)^\top=s(1,0,2)^\top+t(0,1,2)^\top\), with \(s,t\in\R\).
\end{solution}

\begin{problem}\label{ex:06:pythagoras}
In a real inner-product space show that \(x\perp y\) implies
\(\norm{x+y}^2=\norm{x}^2+\norm{y}^2\). Does the converse hold?
\end{problem}
\begin{solution}
Expansion gives \(\norm{x+y}^2-\norm{x}^2-\norm{y}^2=2\ip{x}{y}\).
Thus each condition is equivalent to the other over the real field.
\end{solution}

\begin{problem}\label{ex:06:orthonormal-independent}
Show that the columns of \(I_n\) are orthonormal and independent.
Then prove that every orthonormal list is independent.
\end{problem}
\begin{solution}
For coordinate columns, \(e_i^\top e_j\) equals one for \(i=j\) and zero
otherwise, and \(\sum_i c_i e_i=(c_1,\ldots,c_n)^\top\) vanishes only
when all coefficients vanish. More generally, taking the inner product
of \(\sum_i c_i q_i=0\) with \(q_j\) gives \(c_j=0\) for every \(j\).
\end{solution}

\begin{problem}\label{ex:06:walls}
Two subspaces are orthogonal when every vector of either is orthogonal
to every vector of the other. Explain why a room's front and side walls,
viewed as planes through a corner, are not orthogonal subspaces.
Show also that the zero subspace is orthogonal to every subspace.
\end{problem}
\begin{solution}
A nonzero vector along the common vertical edge belongs to both planes.
It has positive inner product with itself, contradicting subspace
orthogonality. The zero subspace contains only zero, whose inner product
with every vector is zero.
\end{solution}

\begin{problem}\label{ex:06:complement}
For a subspace \(W\) of an inner-product space \(V\), determine
\(\{0\}^\perp\), prove that \(W^\perp\) is a subspace, and show that
\(W\cap W^\perp=\{0\}\).
\end{problem}
\begin{solution}
Every vector is perpendicular to zero, so \(\{0\}^\perp=V\).
If \(x,y\in W^\perp\), then for every \(w\in W\),
\(\ip{ax+by}{w}=a\ip{x}{w}+b\ip{y}{w}=0\); zero also belongs to the
set. Finally a vector in the intersection is perpendicular to itself,
so positive definiteness forces it to be zero.
\end{solution}

\begin{problem}\label{ex:06:orthonormal-basis}
Prove that every finite-dimensional real inner-product space has an
orthonormal basis.
\end{problem}
\begin{solution}
For the zero space use the empty list. Otherwise take a basis
\(v_1,\ldots,v_n\), and recursively set
\(u_j=v_j-\sum_{i<j}\ip{v_j}{u_i}u_i/\norm{u_i}^2\).
Inductively the previously constructed \(u_i\) are nonzero and mutually
orthogonal, and their span equals that of \(v_1,\ldots,v_{j-1}\).
Taking inner products proves \(u_j\perp u_i\). If \(u_j=0\), then
\(v_j\) is in the preceding span, contradicting independence.
The formula and its rearrangement show equality of the two spans at
stage \(j\). Thus all \(u_j\) exist, and \(u_j/\norm{u_j}\) form the
required orthonormal basis.
\end{solution}

\begin{problem}\label{ex:06:gs-three}
Apply Gram--Schmidt, in the given order, to obtain orthonormal bases for
the spans of these lists:
\begin{enumerate}
\item \((1,1,-1)^\top,(1,0,1)^\top\);
\item \((2,1,1)^\top,(1,3,-1)^\top\);
\item \((1,2,1,0)^\top,(1,2,3,1)^\top\);
\item \((1,1,0,0)^\top,(1,-1,1,1)^\top,(-1,0,2,1)^\top\).
\end{enumerate}
\end{problem}
\begin{solution}
(a) The vectors are already orthogonal, with squared norms \(3,2\);
divide them by \(\sqrt3,\sqrt2\).
(b) The first squared norm is \(6\), and the mixed product is \(4\).
The residual is \((1,3,-1)^\top-\tfrac23(2,1,1)^\top
=\tfrac13(-1,7,-5)^\top\). An orthonormal basis is
\((2,1,1)^\top/\sqrt6,\ (-1,7,-5)^\top/\sqrt{75}\).
(c) The mixed product is \(8\), and the first squared norm is \(6\).
Subtracting \(\tfrac43(1,2,1,0)^\top\) leaves
\(\tfrac13(-1,-2,5,3)^\top\). Divide the directions
\((1,2,1,0)^\top,(-1,-2,5,3)^\top\) by \(\sqrt6,\sqrt{39}\).
(d) The first two vectors are orthogonal, with squared norms \(2,4\).
Their inner products with the third are \(-1,2\), respectively, so its
residual is
\[
(-1,0,2,1)^\top+\tfrac12(1,1,0,0)^\top
-\tfrac12(1,-1,1,1)^\top=\tfrac12(-2,2,3,1)^\top.
\]
The required basis consists of the first two vectors divided by
\(\sqrt2,2\), followed by \((-2,2,3,1)^\top/\sqrt{18}\).
\end{solution}

\begin{problem}\label{ex:06:gs-systems}
Find an orthogonal basis for each solution space:
\begin{enumerate}
\item \(2x+y-z=0,\quad y+z=0\);
\item \(x-y+z=0\);
\item \(4x+7y-\pi z=0,\quad2x-y+z=0\);
\item \(x+y+z=0,\quad x-y=0,\quad y+z=0\).
\end{enumerate}
\end{problem}
\begin{solution}
(a) The equations give \((x,y,z)=t(1,-1,1)\); the single displayed
nonzero vector is an orthogonal basis.
(b) Write \((x,y,z)=s(1,1,0)+t(-1,0,1)\).
Subtracting the projection of the second direction onto the first leaves
\((-1/2,1/2,1)\), so \((1,1,0)^\top,(-1,1,2)^\top\) is an
orthogonal basis.
(c) Substitution \(y=2x+z\) gives \(18x+(7-\pi)z=0\).
Thus the entire solution space is spanned by
\((\pi-7,2\pi+4,18)^\top\), which is nonzero.
(d) The second and third equations give \(x=y=-z\); the first forces
\(x=0\). The zero space has the empty orthogonal basis.
\end{solution}

\begin{problem}\label{ex:06:gs-polynomials-two}
For \(\ip{f}{g}=\int_0^1fg\), apply Gram--Schmidt to \(t,t^2\)
and obtain an orthonormal basis of their span.
\end{problem}
\begin{solution}
The first squared norm is \(1/3\), and \(\ip{t^2}{t}=1/4\).
Hence the residual is \(u=t^2-\tfrac34t\), with
\[
\norm{u}^2=\tfrac15-\tfrac32\tfrac14+\tfrac9{16}\tfrac13
=\tfrac1{80}.
\]
Normalize to obtain \(\sqrt3\,t\) and \(\sqrt5(4t^2-3t)\).
The nonzero residual also verifies independence of the original pair.
\end{solution}

\begin{problem}\label{ex:06:gs-polynomials-three}
For \(\ip{f}{g}=\int_0^1fg\), apply Gram--Schmidt to \(1,t,t^2\)
and give an orthonormal basis of their span.
\end{problem}
\begin{solution}
The first vector is \(1\), already normalized. The next residual is
\(u_2=t-\tfrac12\), with squared norm \(1/12\).
For \(t^2\), the products with \(1,u_2\) are \(1/3,1/12\), so
\[
u_3=t^2-\tfrac13-u_2=t^2-t+\tfrac16,
\qquad \norm{u_3}^2
=\tfrac15-\tfrac12+\tfrac49-\tfrac16+\tfrac1{36}=\tfrac1{180}.
\]
Thus an orthonormal basis is
\(1,\ \sqrt3(2t-1),\ \sqrt5(6t^2-6t+1)\).
\end{solution}

\begin{problem}\label{ex:06:gs-dependent}
Let \(A_1,\ldots,A_k\) be matrices in a real inner-product space,
possibly dependent. Show that coefficients can be chosen so that
\(B_j=A_j-\sum_{i<j}c_{ij}B_i\) are mutually orthogonal.
Show that \(c_{ij}=\ip{A_j}{B_i}/\norm{B_i}^2\) is forced when
\(B_i\ne0\), and determine how many \(B_i\) are nonzero.
Describe how to obtain an orthonormal basis for the original span.
\end{problem}
\begin{solution}
Construct successively. For each earlier nonzero \(B_i\), choose the
displayed coefficient, and choose zero when \(B_i=0\).
Orthogonality of the previous list gives
\(\ip{B_j}{B_i}=\ip{A_j}{B_i}-c_{ij}\norm{B_i}^2=0\)
for nonzero \(B_i\); for zero \(B_i\) it is automatic.
The same equation proves uniqueness of each coefficient attached to a
nonzero vector; coefficients attached to zero are arbitrary.
The triangular formulas in both directions show
\(\operatorname{span}(B_1,\ldots,B_j)=\operatorname{span}(A_1,\ldots,A_j)\).
The nonzero \(B_i\) are orthogonal and hence independent. They therefore
form a basis of the original span, and their number is its dimension.
Discard zero vectors and normalize the remaining ones to get the requested
orthonormal basis.
\end{solution}

\begin{problem}\label{ex:06:qr-dependent}
Let \(A\) be a real \(m\times k\) matrix of rank \(r\).
Show that \(A=QR_1\), where \(Q\) has \(r\) orthonormal columns and
\(R_1\) consists of the nonzero rows of an upper triangular \(k\times k\)
matrix \(R\). Require \(R\) to have \(r\) positive diagonal entries and
\(k-r\) zero rows.
\end{problem}
\begin{solution}
Apply the preceding construction to the columns \(a_j\), obtaining
orthogonal residuals \(b_j\). For every nonzero \(b_i\), set
\(q_i=b_i/\norm{b_i}\), keeping its original index for now, and set
\(R_{ij}=\ip{q_i}{a_j}\). For zero \(b_i\), set row \(i\) of \(R\)
equal to zero. A nonzero \(b_i\) is perpendicular to the span of
\(a_1,\ldots,a_{i-1}\), so \(R_{ij}=0\) for \(j<i\);
further \(R_{ii}=\norm{b_i}>0\).
Each \(a_j\) lies in the span of the nonzero residuals through index \(j\),
so its orthogonal expansion is \(a_j=\sum_{i\leqslant j,\ b_i\ne0}q_iR_{ij}\).
Assemble the \(r\) nonzero normalized residuals into \(Q\), and retain
exactly their rows in \(R_1\). These column identities give \(A=QR_1\)
and \(Q^\top Q=I_r\). For \(r=0\), use an empty \(Q,R_1\) and \(R=0\).
\end{solution}

\begin{problem}\label{ex:06:trace-complement}
On the space \(V\) of real symmetric \(n\times n\) matrices, \(n\geqslant1\),
prove that \(\ip{A}{B}=\tr(AB)\) is an inner product.
If \(W=\{A\in V:\tr A=0\}\), prove \(\dim W^\perp=1\).
\end{problem}
\begin{solution}
For symmetric matrices \(\tr(AB)=\tr(A^\top B)=\sum_{i,j}a_{ij}b_{ij}\),
so this is the restriction of the Frobenius inner product.
Also \(\ip{I}{A}=\tr A\), hence \(\operatorname{span}(I)\subseteq W^\perp\).
If \(B\in W^\perp\), let \(C=B-(\tr B/n)I\). Then \(C\in W\), and
\(0=\ip{B}{C}=\ip{C}{C}+(\tr B/n)\tr C=\norm{C}^2\).
Thus \(B=(\tr B/n)I\), so \(W^\perp=\operatorname{span}(I)\), of dimension one.
\end{solution}

\begin{problem}\label{ex:06:dimension}
For a subspace \(W\) of a finite-dimensional inner-product space \(V\),
prove \(\dim W+\dim W^\perp=\dim V\).
\end{problem}
\begin{solution}
Take an orthonormal basis \(q_1,\ldots,q_k\) of \(W\), extend it to a
basis of \(V\), and Gram--Schmidt the added vectors. This leaves the
first \(k\) vectors unchanged and gives an orthonormal basis
\(q_1,\ldots,q_n\) of \(V\). The last \(n-k\) vectors are perpendicular
to \(W\). Conversely, if \(v\perp W\), its coefficients along the first
\(k\) basis vectors vanish, so it is in the span of the last \(n-k\).
They form a basis of \(W^\perp\), proving the formula, including \(k=0,n\).
\end{solution}

\needspace{8\baselineskip}
\subsection*{Forms, matrices, and functions}

\begin{problem}\label{ex:06:positive-form}
For symmetric \(A\), let \(g(x,y)=x^\top Ay\).
\begin{enumerate}
\item Prove symmetry and linearity in each argument.
\item Give a \(1\times1\) and a nonzero \(2\times2\) matrix for which positivity fails.
\item Give a symmetric positive definite \(2\times2\) example.
\item If \(A=\begin{pmatrix}a&b\\b&d\end{pmatrix}\), \(a>0\), prove positivity
is equivalent to \(ad-b^2>0\).
\item If \(a<0\), show that positivity fails.
\end{enumerate}
\end{problem}
\begin{solution}
(a) Matrix multiplication is linear in each vector and
\(x^\top Ay=(x^\top Ay)^\top=y^\top A^\top x=y^\top Ax\).
(b) Take \((-1)\) and \(\operatorname{diag}(1,0)\); the former is negative
on \(1\), and the latter vanishes on nonzero \(e_2\).
(c) \(I_2\) works because \(x_1^2+x_2^2>0\) for nonzero \(x\).
(d) Completing the square gives
\[
ax^2+2bxy+dy^2=a(x+by/a)^2+\frac{ad-b^2}{a}y^2.
\]
If both coefficients are positive, the expression vanishes only at
\((0,0)\). Conversely set \((x,y)=(-b/a,1)\); positive definiteness
forces \((ad-b^2)/a>0\).
(e) Evaluation at \(e_1\) gives \(a<0\), contradicting positivity.
\end{solution}

\begin{problem}\label{ex:06:positive-classification}
Determine which matrices are positive definite:
\[
\begin{array}{lll}
\text{(a) }\begin{pmatrix}3&-1\\-1&2\end{pmatrix}&
\text{(b) }\begin{pmatrix}-2&1\\1&5\end{pmatrix}&
\text{(c) }\begin{pmatrix}4&1\\1&2\end{pmatrix}\\[6pt]
\text{(d) }\begin{pmatrix}4&4\\4&1\end{pmatrix}&
\text{(e) }\begin{pmatrix}4&1\\1&10\end{pmatrix}&
\text{(f) }\begin{pmatrix}4&-1\\-1&10\end{pmatrix}.
\end{array}
\]
\end{problem}
\begin{solution}
In (a), (c), (e), (f), the first diagonal entries are positive and the
values \(ad-b^2\) are respectively \(5,7,39,39\); all four are positive
definite by the preceding square completion. In (b), \(e_1^\top Ae_1=-2\).
In (d), evaluating at \((-1,1)^\top\) gives \(4-8+1=-3\).
Thus (b) and (d) fail positivity.
\end{solution}

\begin{problem}\label{ex:06:symmetric-bilinear}
For a real square matrix \(A\), prove that \(g_A(x,y)=x^\top Ay\) is
symmetric in its arguments if and only if \(A=A^\top\).
\end{problem}
\begin{solution}
If \(A=A^\top\), transposing the scalar \(x^\top Ay\) gives
\(x^\top Ay=y^\top A^\top x=y^\top Ax\).
Conversely, symmetry on \((e_i,e_j)\) gives
\(a_{ij}=g_A(e_i,e_j)=g_A(e_j,e_i)=a_{ji}\) for every \(i,j\).
Thus \(A=A^\top\). This assertion alone imposes no positivity condition.
\end{solution}

\begin{problem}\label{ex:06:bilinear-coordinates}
Write \(x^\top Ay\) fully in coordinates for each matrix:
\[
\begin{array}{lll}
\text{(a) }\begin{pmatrix}2&-3\\4&1\end{pmatrix}&
\text{(b) }\begin{pmatrix}4&1\\-2&5\end{pmatrix}&
\text{(c) }\begin{pmatrix}-5&2\\\pi&7\end{pmatrix}\\[6pt]
\text{(d) }\begin{pmatrix}1&2&-1\\-3&1&4\\2&5&-1\end{pmatrix}&
\text{(e) }\begin{pmatrix}-4&2&1\\3&1&1\\2&5&7\end{pmatrix}&
\text{(f) }\begin{pmatrix}-\frac12&2&-5\\1&\frac23&4\\-1&0&3\end{pmatrix}.
\end{array}
\]
\end{problem}
\begin{solution}
Use \(x^\top Ay=\sum_{i,j}a_{ij}x_i y_j\). The answers are
\begin{align*}
\text{(a)}\quad&2x_1y_1-3x_1y_2+4x_2y_1+x_2y_2,\\
\text{(b)}\quad&4x_1y_1+x_1y_2-2x_2y_1+5x_2y_2,\\
\text{(c)}\quad&-5x_1y_1+2x_1y_2+\pi x_2y_1+7x_2y_2,\\
\text{(d)}\quad&x_1y_1+2x_1y_2-x_1y_3-3x_2y_1+x_2y_2+4x_2y_3\\
&\hspace{1em}+2x_3y_1+5x_3y_2-x_3y_3,\\
\text{(e)}\quad&-4x_1y_1+2x_1y_2+x_1y_3+3x_2y_1+x_2y_2+x_2y_3\\
&\hspace{1em}+2x_3y_1+5x_3y_2+7x_3y_3,\\
\text{(f)}\quad&-\tfrac12x_1y_1+2x_1y_2-5x_1y_3+x_2y_1+\tfrac23x_2y_2+4x_2y_3\\
&\hspace{1em}-x_3y_1+3x_3y_3.
\end{align*}
\end{solution}

\begin{problem}\label{ex:06:trig-orthogonal}
On the continuous real functions on \([0,2\pi]\), use
\(\ip{f}{g}=\int_0^{2\pi}fg\). For \(g_n(t)=\cos(nt)\), \(n\geqslant0\),
and \(h_m(t)=\sin(mt)\), \(m\geqslant1\), compute their norms.
Show that \(g_n\perp g_m\) when \(n\ne m\) and that \(g_n\perp h_m\) always.
\end{problem}
\begin{solution}
For nonzero integer \(k\), the integrals of \(\cos kt\) and \(\sin kt\)
over a full period vanish. Use
\(2\cos nt\cos mt=\cos((n-m)t)+\cos((n+m)t)\) and
\(2\cos nt\sin mt=\sin((m+n)t)+\sin((m-n)t)\).
The first integral is zero for \(n\ne m\), and the second is always zero,
including the zero-frequency term. Also
\(\cos^2 nt=(1+\cos2nt)/2\) for \(n\geqslant1\) and
\(\sin^2 mt=(1-\cos2mt)/2\). Hence
\(\norm{g_0}=\sqrt{2\pi}\), and
\(\norm{g_n}=\norm{h_m}=\sqrt\pi\) for positive indices.
\end{solution}

\begin{problem}\label{ex:06:fourier-linear}
With the preceding product, compute the inner products of \(f(t)=t\)
with \(g_n,h_n\), and its Fourier coefficients
\(\ip{f}{g_n}/\norm{g_n}^2\), \(\ip{f}{h_n}/\norm{h_n}^2\).
\end{problem}
\begin{solution}
For \(n=0\), \(\ip{t}{1}=2\pi^2\), giving coefficient \(\pi\).
For \(n\geqslant1\), integration by parts gives
\[
\int_0^{2\pi}t\cos nt\,dt=
\left[\frac{t\sin nt}{n}+\frac{\cos nt}{n^2}\right]_0^{2\pi}=0,
\]
\[
\int_0^{2\pi}t\sin nt\,dt=
\left[-\frac{t\cos nt}{n}+\frac{\sin nt}{n^2}\right]_0^{2\pi}
=-\frac{2\pi}{n}.
\]
Dividing by the squared norms \(\pi\) gives coefficients \(0,-2/n\).
\end{solution}

\begin{problem}\label{ex:06:fourier-quadratic}
Compute the same inner products and Fourier coefficients for \(f(t)=t^2\).
\end{problem}
\begin{solution}
The constant inner product is \(\int_0^{2\pi}t^2dt=8\pi^3/3\), giving
coefficient \(4\pi^2/3\). For \(n\geqslant1\), integration by parts,
using the two integrals in the preceding solution, gives
\[
\int_0^{2\pi}t^2\cos nt\,dt
=-\frac2n\int_0^{2\pi}t\sin nt\,dt=\frac{4\pi}{n^2},
\]
\[
\int_0^{2\pi}t^2\sin nt\,dt
=-\frac{4\pi^2}{n}+\frac2n\int_0^{2\pi}t\cos nt\,dt
=-\frac{4\pi^2}{n}.
\]
The nonconstant coefficients are therefore \(4/n^2\) and \(-4\pi/n\).
\end{solution}

\begin{problem}\label{ex:06:function-norms}
For the same integral product on \([0,2\pi]\), find the norms of
\(f(t)=t\) and \(g(t)=t^2\).
\end{problem}
\begin{solution}
By definition,
\(\norm{f}^2=\int_0^{2\pi}t^2dt=8\pi^3/3\) and
\(\norm{g}^2=\int_0^{2\pi}t^4dt=32\pi^5/5\).
The norms are the positive square roots of these quantities.
\end{solution}

\begin{problem}\label{ex:06:parseval-basis}
Let \(q_1,\ldots,q_m\) be orthonormal in a finite-dimensional real
inner-product space \(V\). Suppose
\(\norm{v}^2=\sum_i\ip{v}{q_i}^2\) for every \(v\in V\).
Prove that the list is a basis of \(V\).
\end{problem}
\begin{solution}
Orthogonality proves independence. For any \(v\), set
\(p=\sum_i\ip{v}{q_i}q_i\). Direct expansion gives
\[
\norm{v-p}^2=\norm{v}^2-2\sum_i\ip{v}{q_i}^2
+\sum_i\ip{v}{q_i}^2=0.
\]
Thus \(v=p\), proving that the list spans \(V\). It is therefore a basis.
\end{solution}
